\documentclass{article}
% Source: Topic_01_Teacher_Edition.pdf, PDF pages 24-37 (printed pages 13A-22B)
\usepackage{amsmath}
\usepackage{amssymb}
\usepackage{graphicx}
\usepackage{tikz}
\usepackage{pgfplots}
\pgfplotsset{compat=1.18}
\usepackage{booktabs}
\usepackage{xcolor}

\newenvironment{lesson-meta}{}{}        % lesson code, title, objective, EQ, vocab
\newenvironment{item-analysis}{}{}      % Savvas's Example × DOK table
\newenvironment{model-discuss}{}{}      % Launch opener
\newenvironment{concept-box}{}{}        % in-lesson concept reference
\newenvironment{concept-summary}{}{}    % end-of-lesson summary
\newenvironment{example}[1]{}{}         % #1 = example number
\newenvironment{tryit}[1]{}{}           % #1 = tryit number
\newenvironment{practice}[2]{}{}        % #1 = practice number, #2 = DOK
\newenvironment{te-addendum}[2]{}{}     % #1 = anchor example, #2 = type
\newcommand{\answer}[1]{}               % answer key, inside any env
\newcommand{\placeholder}[2]{}          % #1 = type, #2 = description
\newcommand{\uncertain}[2]{#1}          % #1 = best guess, #2 = alternatives
\newcommand{\te}[1]{}                   % teacher-only inline annotation

\begin{document}


% Source page 24: 13A Lesson Overview
\begin{lesson-meta}
lesson: 1-2
title: Transformations of Functions
standards: none printed
practices: none printed
objective: I can apply transformations to graph functions and write equations.

essential question: How do the values a, b, h, and k in (a\cdot f[b(x-h)]+k) transform the graph of a parent function, (f(x))?

\textbf{VOCABULARY}
\item compression
\item reflection
\item stretch
\item transformation
\item translation
\textbf{TOPIC VOCABULARY}
\te{No topic vocabulary list is printed on these pages. The I CAN statement is on PDF page 25 and Essential Question on page 26. No numbered standards or Mathematical Practice codes are printed.}
\begin{tabular}{ll}
Lesson Code & 1-2 \\
Title & Transformations of Functions \\
Standards & none printed \\
I CAN Objective & I can apply transformations to graph functions and write equations. \\
Essential Question & How do the values a, b, h, and k in (a\cdot f[b(x-h)]+k) transform the graph of a parent function, (f(x))? \\
Lesson Vocabulary & compression; reflection; stretch; transformation; translation \\
Topic Vocabulary & none printed \\
\end{tabular}
\end{lesson-meta}
\begin{te-addendum}{0}{GeneralizeETP}
\textbf{Lesson Overview: FOCUS}
\textbf{Objective} Students will be able to:
\item Graph a transformed function by identifying the effect on the graph of replacing (f(x)) by (f(x)+k), (kf(x)), (f(kx)), and (f(x+k)) for specific values of (k).
\item Write an equation of a transformed function.
\item Relate the domain of a function to its graph and the real-world situation it describes.
\textbf{Essential Understanding}
A function of the form (f(x)=a\cdot f[b(x-h)]+k) is transformed by changing the values of (a), (b), (h), or (k). Changing the value of (h) or (k) results in a horizontal or vertical translation, changing the sign of (a) or (b) results in a reflection across one of the axes, and changing the value of (a) or (b) results in a horizontal or vertical stretch or compression.
\textbf{COHERENCE}
Previously in earlier courses, students:
\item Created, graphed, and identified the key features of the graphs of linear, quadratic, and absolute value functions.
In this lesson, students:
\item Write equations for linear and quadratic functions by understanding how changing the values of (a), (b), (h), and (k) affect the key features of the graph of a function.
Later in this topic, students will:
\item Create, graph, and identify the key features of piecewise-defined functions such as absolute value functions and step functions.
\textbf{RIGOR}
This lesson emphasizes a blend of conceptual understanding and application.
\item Students apply their understanding of transformations either to graph a translation of a linear, quadratic, or absolute value function or to write an equation given a graph.
\end{te-addendum}
\begin{te-addendum}{0}{ELLAddendum}
\textbf{Vocabulary Builder}
Glossary is available at SavvasRealize.com.
\textbf{NEW VOCABULARY}
\begin{tabular}{ll}
English & Spanish \\
compression & compresión \\
reflection & reflexión \\
stretch & tramo \\
transformation & transformación \\
translation & translación \\
\end{tabular}
\textbf{VOCABULARY ACTIVITY}
Using a piece of clay or similar material, provide a visual example of the terms compression and stretch. Tell students that a transformation of a function in the form (f(x)=a\cdot f[b(x-h)]+k) involves a change to at least one of the values (a), (b), (h), or (k). Tell them that some transformations affect the vertical distance between the points of the graph of a function. When the points are closer together, it is called a compression, and when the points are further apart, it is called a stretch. Have students complete the following sentences.
A transformation that decreases the distance between the points of the graph is called a blank.
\answer{compression}
A transformation that increases the distance between the points of the graph is called a blank.
\answer{stretch}
\end{te-addendum}
\begin{te-addendum}{0}{LearnTogether}
\textbf{Student Companion}
Students can do their work for the lesson in their Student Companion or on Savvas Realize.
\placeholder{illustration}{Student Companion cover, dark background with colored loops, white title Student Companion and enVision Algebra 2.}
\end{te-addendum}
\begin{te-addendum}{0}{GeneralizeETP}
\textbf{Mathematics Overview}
Students interpret, graph, and write equations of transformations of functions. Functions in the form (y=f(x-h)) represent horizontal translations, functions in the form (y=f(x)+k) represent vertical translation, functions in the form (y=-f(x)) represent reflections across the x-axis, functions in the form (y=f(-x)) represent reflections across the y-axis, functions in the form (y=f(bx)) represent horizontal stretches or compressions, and functions in the form (y=af(x)) represent vertical stretches or compressions.
\textbf{Applying Math Practices}
\textbf{Use Appropriate Tools Strategically}
Students use available tools, such as graphing calculators, to graph the original and transformed equations to check that the transformations they identified are correct.
\textbf{Look For and Make Use of Structure}
Students look for overall structure and patterns in mathematics when identifying why the transformation (g(x)=-f(x)) affects the y-coordinate of each point and why the transformation (g(x)=f(-x)) affects the x-coordinate of each point.
\te{SavvasRealize.com: Program Resources.}
\end{te-addendum}

% Source page 25: 13B STEP 1 Explore
\begin{te-addendum}{0}{LearnTogether}
\textbf{EXPLORE \& REASON}
\textbf{INSTRUCTIONAL FOCUS} Students explore the graph of an absolute value function and how different values for (c) impact the graph of the function.
\textbf{STUDENT COMPANION} Students can complete the Explore \& Reason activity in their Student Companion.
\end{te-addendum}
\begin{te-addendum}{0}{GeneralizeETP}
\textbf{Before—WHOLE CLASS}
\textbf{Implement Tasks that Promote Reasoning and Problem Solving ETP}
Q: What does the graph of the function tell you about the range of (f(x)=|x|)?
\answer{The range is greater than or equal to zero.}
Q: What can you predict about the range of the graphs of the function (g(x)=|x+c|) for different values of (c)?
\answer{Regardless of the value of (c), the range will still be greater than or equal to zero.}
\textbf{During—SMALL GROUP}
\textbf{Support Productive Struggle in Learning Mathematics ETP}
Q: How can you find points to graph for each function?
\answer{Make a table of values using values of (x) between (-5) and 5 and calculate the corresponding y-values.}
Q: What do you notice about the vertex of each function?
\answer{The vertices all lie on the x-axis.}
\textbf{For Early Finishers}
Q: How does the graph change if the coefficient of (f(x+c)) is negative? Choose three of your functions and graph (-f(x+c)).
\answer{The graphs of (-f(x+c)) are reflections across the x-axis of the graphs of (f(x+c)).}
\textbf{After—WHOLE CLASS}
\textbf{Facilitate Meaningful Mathematical Discourse ETP}
Discuss how transformations applied to functions can affect the range, the domain, or both.
Q: If the function (f(x)) has a horizontal shift, what part of the function has been changed?
\answer{In the function (f(x)), (x) stands for the domain; therefore, operations will be applied to (x) within the parentheses.}
Q: If the function (f(x)) has a vertical shift, what part of the function has been changed?
\answer{In the function (f(x)), (f(x)) stands for the range; therefore, operations will be applied to (f(x)).}
Introduce students to the language of functions written in the form (f(x)=a\cdot f[b(x-h)]+k) and describe how that notation is used to define certain movements of the graph of a function.
\end{te-addendum}
\begin{te-addendum}{0}{HabitsOfMind}
\textbf{Use with EXPLORE \& REASON}
Q: Reason What happens to the graph if (c) is a negative number?
\answer{The graph of the function is shifted to the right when (c) is a negative number.}
\end{te-addendum}
\begin{model-discuss}
\textbf{EXPLORE \& REASON}
The graph of the function (f(x)=|x|) is shown.
\placeholder{graph}{Student Edition launch: red V f(x)=|x|, vertex (0,0), slopes -1 and 1, upper ends arrowed. x/y axes arrowed with origin O and labels -1 and 1; unit square grid, visible window approximately [-2,2] by [-2,2]. No point circles.}
A. Graph the function (g(x)=|x+c|) several times with different values for (c) (any value from (-5) to 5).
\answer{A. graphs with several absolute value functions identical to (f(x)=|x|) but shifted several units to the left or right}
B. Look for Relationships Predict what will happen to the graph if (c) is a number greater than 100. What if (c) is a number between 0 and (\frac{1}{2})?
\answer{B. For (c>100), you will see a very large horizontal shift of the graph to the left. For (0<c<\frac{1}{2}) you will see a very small horizontal shift of the graph to the left.}
\te{Printed SAMPLE STUDENT WORK supplies the answers. Explore \& Reason is Powered by Desmos. Activities are available at SavvasRealize.com. The digital copy repeats the prompt with “Use the tool to graph” in part A and an “Enter your answer.” field.}
\placeholder{graph}{Digital launch screenshot: blue V f(x)=|x| with vertex (0,0); unit grid and major lines at multiples of 5, x labels -10,-5,0,5,10 and y labels -5,5. Window approximately x=-10 to 10, y=-8 to 8. No point circles or curve arrowheads. Desmos controls at upper right and Go back label above.}
\end{model-discuss}


% Source page 26: 13 STEP 2 Understand & Apply
\begin{te-addendum}{0}{LearnTogether}
\textbf{INTRODUCE THE ESSENTIAL QUESTION}
\textbf{Establish Mathematics Goals to Focus Learning ETP}
Introduce students to the different types of transformations of functions. Remind them that transformations are defined by vertical and horizontal translations, reflections, and stretches and compressions on the coordinate plane. Students can apply these transformations to graphing functions and writing equations.
\textbf{ESSENTIAL QUESTION} How do the values (a), (b), (h), and (k) in (a\cdot f[b(x-h)]+k) transform the graph of a parent function, (f(x))?
\te{Student-page vocabulary repeats compression, reflection, stretch, transformation, translation. Example 1 is Powered by Desmos. Examples are available at SavvasRealize.com.}
\end{te-addendum}
\begin{te-addendum}{1}{GeneralizeETP}
\textbf{EXAMPLE 1 Translate a Function}
\textbf{Use and Connect Mathematical Representations ETP}
Q: What do you notice about the two new graphs compared to their parent functions?
\answer{One graph is the parent function translated down; the other graph is the parent function translated to the right.}
Q: Will the transformation of either graph cause the parabola to change shape?
\answer{No; a translation is a slide to a new position. The shape will not change.}
\end{te-addendum}
\begin{te-addendum}{3}{PurposefulQuestions}
\textbf{Additional Examples—ADDITIONAL EXAMPLE 3}
Graph the function (f(x)=x). Then, use a table to graph (g(x)=f(4x)) and (h(x)=f(\frac{1}{4}x)). Identify the relationship of the horizontal compression and stretch to the parent function.
\begin{tabular}{llllll}
(x) & (f(x)) & (x) & (g(x)) & (x) & (h(x)) \\
(-4) & (-4) & (-1) & (-4) & (-16) & (-4) \\
(-3) & (-3) & & (-3) & (-12) & (-3) \\
(-2) & (-2) & & (-2) & (-8) & (-2) \\
(-1) & (-1) & & (-1) & (-4) & (-1) \\
0 & 0 & 0 & 0 & 0 & 0 \\
1 & 1 & & 1 & 4 & 1 \\
2 & 2 & & 2 & 8 & 2 \\
3 & 3 & & 3 & 12 & 3 \\
4 & 4 & 1 & 4 & 16 & 4 \\
\end{tabular}
\placeholder{graph}{Additional Example 3: green line f(x)=x, orange steep line g(x)=4x, blue shallow line h(x)=x/4, all through origin with arrows both ends and color-matched labels. Square grid spacing 2, axes x/y arrowed, labels -8,4,8 on x and -8,4,8 on y, origin O; window approximately [-10,10] by [-10,10]. No point circles.}
\answer{The graph of (g(x)) is a horizontal compression by a factor of 4. (h(x)) is a horizontal stretch by a factor of 4.}
\textbf{Example 3} Help students understand stretches and compressions of linear functions.
Q: What do you notice about the output values of (g(x)) and (h(x)) compared with the output values of (f(x))?
\answer{For the same input value, the output value of (g) is (\frac{1}{3}) the output value of (f) and the output value of (h) is 3 times the output value of (f).}
\te{Printed factors 1/3 and 3 in the teacher answer conflict with the displayed functions; likely factors 4 and 1/4 for g and h, respectively. Blank x cells in the displayed g table are retained.}
\end{te-addendum}
\begin{te-addendum}{4}{PurposefulQuestions}
\textbf{ADDITIONAL EXAMPLE 4}
What transformations of the graph of (f(x)) result in the graph of the function (p)?
(p(x)=f(\frac{1}{3}x-2)+7)
\placeholder{graph}{Additional Example 4 screenshot: black polyline f(x) joins filled points (-6,2), (-5,0), (-2,4), (0,5). Unit grid, horizontal labels -5,0,5,10,15,20; vertical labels 5,10,15. Window approximately x=-7 to 20, y=-2 to 16. No curve arrows; axes have no visible arrowheads. SOLUTION button at left; no worked solution displayed.}
\te{No answer is displayed in this screenshot.}
\textbf{Example 4} Help students understand functions that have more than two transformations. This example shows how to organize an approach to graphing more complex functions.
Q: How can you determine the horizontal transformations of the graph of (f)?
\answer{In the function, the coefficient of (\frac{1}{3}) indicates a horizontal stretch. Factor (\frac{1}{3}) to find the number subtracted from (x), which indicates the horizontal translation.}
\te{Both screenshots contain Go back, ESSENTIAL QUESTION, and SOLUTION. Additional Examples are available at SavvasRealize.com.}
\end{te-addendum}
\begin{example}{1}
\textbf{Translate a Function}
\textbf{CONCEPTUAL UNDERSTANDING}
A. Graph the function (f(x)=x^2) for the domain ([-2,2]). The graph of (g) is the graph of (f) after a translation of 3 units down. How are the functions, domains, and ranges of (f) and (g) related?
Every point on the graph of (g) is 3 units below a corresponding point on the graph of (f).
(g(x)=f(x)-3)
(g(x)=x^2-3)
\placeholder{graph}{Example 1A: blue parabola f(x)=x^2 from filled (-2,4) to filled (2,4), vertex (0,0); red parabola g(x)=x^2-3 from filled (-2,1) to filled (2,1), vertex (0,-3). Filled corresponding points (1,1) and (1,-2) labeled (x,f(x)) and (x,g(x)). Unit square grid, x/y axes arrowed; labels -4,-2,2,4 on x and -4,2,4 on y, origin O; window approximately [-5,5] by [-5,5].}
\te{Callout: The difference between (f(x)) and (g(x)) is 3.}
The domains of (f) and (g) are the same: ([-2,2]). The range values of (g) are 3 units less than the range values of (f). The range of (f) is ([0,4]) and the range of (g) is ([-3,1]).
A translation is a type of transformation of a function, one that shifts each point on a graph the same distance and direction.
\te{Callout: Other types of transformations include reflections, stretches, and compressions.}
In general, if (g(x)=f(x)+k), then the graph of (g) is a vertical translation of the graph of (f) by (k) units.
\te{LOOK FOR RELATIONSHIPS: This type of transformation slides the graph up or down. You could perform a transformation like this to a non-vertical line and produce a parallel line.}
\answer{A. (g(x)=x^2-3); both domains ([-2,2]); range of (f): ([0,4]); range of (g): ([-3,1]).}
% Source page 27: 14 Example 1 continued
B. Graph the function (f(x)=x^2) for the domain ([-2,2]). The graph of the function (g) is the graph of (f) after a translation 3 units to the right. How are the functions, domains, and ranges of (f) and (g) related?
Every point on the graph of (g) is 3 units to the right of the corresponding point on the graph of (f). Express (g(x)) in terms of (f(x)).
(g(x)=f(x-3))
(g(x)=(x-3)^2)
\placeholder{graph}{Example 1B: blue f(x)=x^2 on [-2,2], filled endpoints (-2,4),(2,4), vertex (0,0); green g(x)=(x-3)^2 on [1,5], filled endpoints (1,4),(5,4), vertex (3,0). Filled points (1,1) and (4,1) labeled (x,f(x)) and (x,g(x)) joined by horizontal dashed segment. Unit grid; x/y axes arrowed, origin O, x labels -4,-2,2,4,6 and y labels -4,-2,2,4. Window approximately x=-5 to 7,y=-5 to 5.}
\te{Callout: The difference between x-coordinates is 3, but the y-coordinates are the same.}
The range values of (f) and (g) are the same: ([0,4]). The domain values of (g) are 3 units more than the domain values of (f). The domain of (f) is ([-2,2]) and the domain of (g) is ([1,5]).
In general, if (g(x)=f(x-h)), then the graph of (g) is a horizontal translation of the graph of (f) by (h) units.
\answer{B. (g(x)=(x-3)^2); both ranges ([0,4]); domain of (f): ([-2,2]); domain of (g): ([1,5]).}
\end{example}
\begin{tryit}{1}
a. Compare the intercepts of (f) to the intercepts of (g) in part (a).
b. Compare the intercepts of (f) to the intercepts of (g) in part (b).
\answer{a. From (f) to (g), the y-intercept is translated down 3 units. (f) had one x-intercept at the origin, and (g) has two x-intercepts on either side of the origin. b. From (f) to (g), the x-intercept is translated right 3 units. (f) had one y-intercept at the origin, but (g) does not have a y-intercept.}
\end{tryit}
\begin{te-addendum}{1}{ElicitEvidence}
\textbf{Elicit and Use Evidence of Student Thinking ETP}
Q: Will a horizontal translation always result in the translated function not having a y-intercept?
\answer{No; if the domain of (f) were not restricted, then any horizontal translation of (f) would still have a y-intercept.}
\end{te-addendum}
\begin{te-addendum}{2}{PurposefulQuestions}
\textbf{EXAMPLE 2 Reflect a Function Across the x- or y-Axis}
\textbf{Pose Purposeful Questions ETP—PART A}
Q: How can you describe what the reflection across the x-axis does to the coordinates and the equation?
\answer{The y-values of the original function are multiplied by (-1).}
\end{te-addendum}
\begin{te-addendum}{2}{ELLAddendum}
\textbf{English Language Learners (Use with EXAMPLE 2)}
\textbf{SPEAKING—BEGINNING} Ask students to tell what a reflection is. Then give each student a mirror and have them draw the reflections of different lines and shapes on a coordinate grid—some across the x-axis and some across the y-axis. Have them describe their reflections.
Q: How do the points of a graph change when reflected across the x-axis? The y-axis?
\answer{The y-coordinates change to the opposite sign when reflected across the x-axis. The x-coordinates change to the opposite sign when reflected across the y-axis.}
Q: How are reflections related to symmetry?
\answer{Symmetrical shapes can be divided by a line down the middle, with one side a reflection of the other side.}
\textbf{LISTENING—DEVELOPING} Read the second paragraph of Part A aloud as students listen. Then play a game where students listen to directions about sending game pieces to different positions on a coordinate grid. For example, “Place the game piece on the point (4,2). Now, send the piece to the opposite y-value.”
Q: When you move from a positive y-value to a negative y-value, in which direction do you move?
\answer{down}
Q: When you move from a negative x-value to a positive x-value, in which direction do you move?
\answer{right}
\textbf{WRITING—EXPANDING} Sometimes words are spelled the same but are pronounced differently and have different meanings. Point out the word produce in the example. Have students write two sentences using the word in two different ways.
Q: What topic could you write about if you used the word produce pronounced with the stress on the first syllable?
\answer{fruits and vegetables; grocery store}
Q: What is a synonym for produce as it is used in the example?
\answer{make, create, form}
\te{Example 2 is Powered by Desmos. Examples are available at SavvasRealize.com.}
\end{te-addendum}
\begin{example}{2}
\textbf{Reflect a Function Across the x- or y-Axis}
A. Graph (f(x)=2x-6) for the domain ([0,\infty)) and the function (g), whose graph is the reflection of the graph of (f) across the x-axis. How are their functions related?
Graph (f). Then graph (g) by reflecting each point of the graph of (f) across the x-axis. For each point ((x,y)) on the graph of (f), plot the point ((x,-y)) to get the graph of (g).
(g(x)=-f(x))
(\phantom{g(x)}=-(2x-6))
(g(x)=-2x+6)
\te{Callout: The y-values of (g) have the opposite sign of (f).}
\placeholder{graph}{Example 2A: red ray f(x)=2x-6 starting at filled (0,-6) and extending up-right; blue g(x)=-2x+6 starting at filled (0,6) and extending down-right. Green filled points (0,6),(0,-6),(3,0),(5,4),(5,-4), last two labeled (x,y),(x,-y) connected by dashed vertical green line. Horizontal reflection axis highlighted green dashed. Unit square grid; x/y axes arrowed, x labels -4,-2,2,4, y labels -4,-2,2,4, origin O; window approximately [-6,6] by [-6,6].}
\te{Callout: The y-intercept of (g), 6, is the opposite of the y-intercept of (f), (-6).}
\te{Callout: The slope of the graph of (f) is 2, while the slope of the graph of (g) is (-2).}
The expression that defines (g) is the opposite of the expression that defines (f). The domain of (g) is the same as (f). The range of (f) is ([-6,\infty)) and the range of (g) is ((-\infty,6]).
\te{VOCABULARY: Recall that a reflection is a transformation that maps each point to a new point across a given line, called the line of reflection.}
\answer{A. (g(x)=-2x+6); same domain as (f); range of (f): ([-6,\infty)); range of (g): ((-\infty,6]).}
% Source page 28: 15 Example 2 continued
B. Graph (f(x)=2x-6) for the domain ([0,\infty)) and the function (h), whose graph is the reflection of the graph of (f) across the y-axis. How are their functions related?
Graph (f). Then reflect every point on the graph of (f) over the y-axis to produce the graph of (h).
(h(x)=f(-x))
(\phantom{h(x)}=2(-x)-6)
(h(x)=-2x-6)
The function (h) has a slope that is the opposite of the slope of (f) but with the same y-intercept. The domain of (h) is ((-\infty,0]). The range for (f) and (h) are the same, ([-6,\infty)).
\placeholder{graph}{Example 2B: red ray f(x)=2x-6 from filled (0,-6) to upper right, blue h(x)=-2x-6 from same point to upper left. Green filled (-3,0),(3,0),(0,-6), and corresponding (-2.5,-1),(2.5,-1) labeled (-x,y),(x,y) with dashed horizontal guide. Vertical reflection axis highlighted green. Unit grid, x/y axes arrowed and origin O, labels -4,-2,2,4 on x, 2,4 on y; window approximately [-6,6] by [-6,6].}
\te{Callout: The slope of the graph of (f) is 2, and the slope of the graph of (h) is (-2).}
\te{Callout: For any point ((x,y)) on the graph of (f), there is a reflected point ((-x,y)) on the graph of (h), so (h(x)=f(-x)). The x-intercept of (h), (-3), is the opposite of the x-intercept of (f), 3.}
\te{USE STRUCTURE: Why does (g(x)=-f(x)) affect the y-coordinate of each point and (g(x)=f(-x)) affect the x-coordinate of each point?}
\answer{B. (h(x)=-2x-6); domain of (h): ((-\infty,0]); both ranges ([-6,\infty)).}
\end{example}
\begin{tryit}{2}
What is the function for the reflected graph? Check by graphing.
a. the graph of (f(x)=x^2-2) reflected across the x-axis
b. the graph of (f(x)=x^2-2) reflected across the y-axis
\answer{a. (y=-x^2+2). b. (y=x^2-2).}
\end{tryit}
\begin{te-addendum}{2}{PurposefulQuestions}
\textbf{PART B}
Q: How can you describe what the reflection across the y-axis does to the coordinates and the equation?
\answer{The x-values of the original function are multiplied by (-1).}
Q: What is the difference between (g(x)=-f(x)) and (g(x)=f(-x))?
\answer{When (f(x)) has a coefficient of (-1), the y-values change sign, causing a reflection across the x-axis. When (x) has a coefficient of (-1), the x-values change sign, causing a reflection across the y-axis.}
\end{te-addendum}
\begin{te-addendum}{2}{CommonError}
\textbf{Common Error}
\textbf{Try It! 2b} Some students may incorrectly interpret a reflection across the y-axis as the result of changing the coefficient of (x) to (-1). Show students that the (x) in the function is not changed to a negative; the actual values that are substituted into the function will be opposites of those from the original function but will produce the same range.
\end{te-addendum}
\begin{te-addendum}{2}{HabitsOfMind}
\textbf{Use with EXAMPLES 1 \& 2}
Q: Look for Relationships How are the intercepts of a graph affected by reflection across the x-axis? Explain.
\answer{The x-intercept stays the same because the points on the axis you are reflecting about do not change. The y-intercept of the reflected graph is the opposite of the y-intercept of the original graph because the two points are equidistant from the x-axis.}
\end{te-addendum}
\begin{te-addendum}{3}{GeneralizeETP}
\textbf{EXAMPLE 3 Understand Stretches and Compressions}
\textbf{Build Procedural Fluency From Conceptual Understanding ETP—PART A}
Q: What is the relationship between the coefficient of (f(x)=a|x|) and its effects on the graph?
\answer{The coefficient of (f(x)) is the factor by which the graph is either stretched or compressed vertically with regard to each point’s distance from the x-axis.}
\end{te-addendum}
\begin{example}{3}
\textbf{Understand Stretches and Compressions}
A. Graph (f(x)=|x|) with domain ([-4,4]) and (g(x)=2\cdot f(x)). How are the graphs, domains, and ranges related?
Use a table to find points on the graph of (g).
\begin{tabular}{llll}
(x) & (f(x)) & (g(x)=2\cdot f(x)) & ((x,g(x))) \\
(-4) & 4 & (2\cdot f(-4)=2(4)=8) & (-4,8) \\
(-2) & 2 & (2\cdot f(-2)=2(2)=4) & (-2,4) \\
0 & 0 & (2\cdot f(0)=2(0)=0) & (0,0) \\
2 & 2 & (2\cdot f(2)=2(2)=4) & (2,4) \\
4 & 4 & (2\cdot f(4)=2(4)=8) & (4,8) \\
\end{tabular}
\placeholder{graph}{Example 3A: blue V f(x)=|x| with filled endpoints (-4,4),(4,4), red V g(x)=2|x| with filled endpoints (-4,8),(4,8), common vertex (0,0). Unit grid; x/y axes arrowed with origin O, x labels -4,-2,2,4, y labels 4,6,8. Window approximately x=-5 to 5,y=-1 to 9; no curve arrows.}
A transformation that increases the distance between the points of a graph and a given line, usually the x-axis, by the same factor is called a stretch. The graph of (g) is a vertical stretch of the graph of (f) by a factor of 2.
The domains of (f) and (g) are the same. The range of (f), ([0,4]), is doubled for (g) to ([0,8]).
\answer{A. Vertical stretch by a factor of 2; same domains; range of (f): ([0,4]); range of (g): ([0,8]).}
% Source page 29: 16 Example 3 continued
B. Graph (f(x)=|x|) with domain ([-4,4]) and (h(x)=f(2x)). How are the graphs, domains, and ranges related?
Use a table to find points on the graph of (h).
\begin{tabular}{llll}
(x) & (f(x)) & (h=f(2x)) & ((x,h(x))) \\
(-2) & 2 & (f(2(-2))=f(-4)=4) & (-2,4) \\
(-1) & 1 & (f(2(-1))=f(-2)=2) & (-1,2) \\
0 & 0 & (f(2(0))=f(0)=0) & (0,0) \\
1 & 1 & (f(2(1))=f(2)=2) & (1,2) \\
2 & 2 & (f(2(2))=f(4)=4) & (2,4) \\
\end{tabular}
\placeholder{graph}{Example 3B: blue V f(x)=|x| with filled endpoints (-4,4),(4,4), green V h(x)=|2x| with filled endpoints (-2,4),(2,4), common vertex (0,0). Unit grid, x/y axes arrowed, origin O, x labels -4,-2,2,4 and y labels 4,6,8; window approximately x=-5 to 5,y=-1 to 9; no curve arrows.}
A transformation that decreases the distance between the points of a graph and a given line, usually the y-axis, by the same factor is called a compression. The graph of (h) is a horizontal compression of the graph of (f) by a factor of 2.
For each corresponding output, the value of the input for (h) is half the value of the input for (f). The two functions have the same range. The values in the domain of (h), ([-2,2]), are half as large as the values in the domain of (f), ([-4,4]).
\te{COMMON ERROR: Be careful not to assume that the domain of a transformed function is the same as the domain of the original function. Notice that (h(-4)) would equal (f(-8)), which is outside the domain of (f), ([-4,4]).}
\answer{B. Horizontal compression by a factor of 2; same range; domain of (h): ([-2,2]); domain of (f): ([-4,4]).}
\end{example}
\begin{tryit}{3}
The domain of (j(x)) is ([a,b]). Show that (j(x)=f(\frac{1}{2}x)) is a horizontal stretch of the graph of (f). What is the domain of (f)?
\answer{If ((x,y)) is on the graph of (j), then ((\frac{1}{2}x,y)) is on the graph of (f). So, every point on (f) corresponds to a point on (j) whose y-coordinate is the same and whose x-coordinate is doubled. So, the graph of (j) is a horizontal stretch of (f). If ([a,b]) is the domain of (j), then ([\frac{1}{2}a,\frac{1}{2}b]) is the domain of (f).}
\end{tryit}
\begin{te-addendum}{3}{GeneralizeETP}
\textbf{PART B}
Q: What is the relationship between the coefficient in the function (f(x)=|bx|) and the effects on the graph of the function?
\answer{When (b>1), the graph is compressed horizontally by a factor (\frac{1}{b}). When (b<1), the graph is stretched by (\frac{1}{b}).}
\te{Printed “When b<1” without a lower bound; likely intended (0<b<1), as in the concept box. The page uses reciprocal-factor wording here and “factor of 2” in Example 3B; both are retained.}
\end{te-addendum}
\begin{concept-box}
\textbf{CONCEPT Stretches and Compressions}
\begin{tabular}{ll}
Vertical & Horizontal \\
(g(x)=k\cdot f(x)) & (h(x)=f(kx)) \\
stretch when (k>1) & stretch when (0<k<1) \\
compression when (0<k<1) & compression when (k>1) \\
\end{tabular}
\placeholder{graph}{Vertical stretch/compression schematic: black cubic-like curve crosses x-axis at approximately -2,0,2 grid units, with local maximum left of origin and minimum right. Red curve has same zeros but taller peak and deeper trough; blue curve same zeros with flatter peak/trough. All curve ends arrowed, lower left and upper right. Square grid and arrowed x/y axes with origin O, no numbered ticks; window about ten grid squares in each direction. No points labeled.}
\placeholder{graph}{Horizontal stretch/compression schematic: black cubic-like curve through origin, local maximum left and minimum right. Red stretched version spreads turning points and other zeros farther horizontally from y-axis; blue compressed version brings them closer. Red visible curve approaches x-axis at left and right borders; black/blue tails extend down and up with arrows. Square grid, arrowed x/y axes and origin O, no numbered ticks or marked points.}
\end{concept-box}
\begin{te-addendum}{3}{RtISupport}
\textbf{Support Struggling Students}
\textbf{USE WITH EXAMPLE 3} For extra support, have students practice creating a table to find points of the function when the domain of (f) is ([-5,5]).
\item Use the function (h(x)=\frac{1}{2}f(x)) to build a table and describe the range and domain of the function.
Q: Describe the domain and range of (h).
\answer{The domain of (h) is the same as that of (f), and the range is (\frac{1}{2}) the range of (f).}
Q: How will the transformation affect the points on the graph of (h) when compared to the graph of (f)?
\answer{Each corresponding point on the graph of (h) is half the distance from the x-axis as the points on the graph of (f).}
\te{Examples are available at SavvasRealize.com.}
\end{te-addendum}


% Source page 30: 17 STEP 2 Understand & Apply
\begin{te-addendum}{4}{PurposefulQuestions}
\textbf{EXAMPLE 4 Graph a Combination of Transformations}
\textbf{Pose Purposeful Questions ETP—PART A}
Q: According to the equation, how many transformations are carried out?
\answer{There are two transformations: a horizontal translation two units left, indicated by the (+2) to (x); and a reflection over the x-axis, indicated by the (-1) as the coefficient of (f(x)).}
Q: How is the reflection created after the translation has occurred?
\answer{Each point on the reflected image is exactly the same distance below the x-axis as the corresponding point of the original image is above the x-axis.}
\textbf{PART B}
Q: For the vertical compression, what distance is affected?
\answer{Each point of the original function will be (\frac{1}{2}) the distance from the x-axis for the vertically compressed image.}
\end{te-addendum}
\begin{example}{4}
\textbf{Graph a Combination of Transformations}
Using the graph (y=f(x)), how can you graph a combination of transformations?
\placeholder{graph}{Original blue polyline f joins (-4,2),(-2,4),(-1,3),(1,7), no marked endpoints or arrows on line. Unit grid with x/y axes arrowed and origin O; x labels -6,-4,-2,2,4,6; y labels -6,-4,-2,2,6; window approximately [-7,7] by [-7,7].}
A. Graph (y=-f(x+2)).
Graph (g(x)=f(x+2)), which is a translation of (f) left 2 units.
Graph (h(x)=-f(x+2)), which is a reflection of (g) across the x-axis.
The graph of (h) is a translation of (f) left 2 units followed by a reflection across the x-axis.
\placeholder{graph}{Example 4A: same unit-grid axes/window as original. Blue f polyline (-4,2),(-2,4),(-1,3),(1,7). Green dashed g polyline (-6,2),(-4,4),(-3,3),(-1,7). Red h polyline (-6,-2),(-4,-4),(-3,-3),(-1,-7). Color labels f,g,h. No endpoint circles or curve arrows.}
\answer{A. Translate (f) left 2 units, then reflect across the x-axis.}
B. Graph (y=\frac{1}{2}f(x)+3).
Graph (g(x)=\frac{1}{2}f(x)), which is a vertical compression of (f) by the factor (\frac{1}{2}).
Graph (h(x)=\frac{1}{2}f(x)+3), which is a translation of (g) up 3 units.
The graph of (h) is a vertical compression of (f) by the factor (\frac{1}{2}) followed by a translation 3 units up.
\placeholder{graph}{Example 4B: same unit-grid axes/window as original. Blue f polyline (-4,2),(-2,4),(-1,3),(1,7). Green dashed g polyline (-4,1),(-2,2),(-1,1.5),(1,3.5). Red h polyline (-4,4),(-2,5),(-1,4.5),(1,6.5). Color labels f,g,h. No endpoint circles or curve arrows.}
\answer{B. Vertically compress (f) by factor (\frac{1}{2}), then translate up 3 units.}
\te{STUDY TIP: To graph transformations with multiple steps, graph the intermediate graphs following equation order of operations. Parentheses first, then exponents, multiplication and division from left to right, and finally addition and subtraction from left to right.}
\end{example}
\begin{tryit}{4}
Using the graph of (f) above, graph each equation.
a. (y=f(2x)-4)
b. (y=f(2x-3)-2)
\answer{a. See graph. b. See graph.}
\placeholder{graph}{Try It 4a answer: blue original f polyline (-4,2),(-2,4),(-1,3),(1,7), red g polyline (-2,-2),(-1,0),(-0.5,-1),(0.5,3). Unit square grid with arrowed x/y axes and origin O, x labels -4,-2,2,4, y labels -2,2,6; window approximately x=-5 to 5,y=-3 to 7. No endpoint circles or curve arrows.}
\placeholder{graph}{Try It 4b answer: same blue original f and same grid/window as 4a. Red g polyline (-0.5,0),(0.5,2),(1,1),(2,5). Axes x/y arrowed, origin O, x labels -4,-2,2,4, y labels -2,2,6. No endpoint circles or curve arrows.}
\end{tryit}
\begin{te-addendum}{4}{ElicitEvidence}
\textbf{Elicit and Use Evidence of Student Thinking ETP}
Q: In part (b) what order would be best for carrying out the transformations?
\answer{(f(2x-3)-2=f(2(x-1.5))-2), so apply the horizontal compression by a factor of 2 first, then the horizontal shift right 1.5 units, and last, the vertical shift down 2 units.}
\end{te-addendum}
\begin{te-addendum}{4}{HabitsOfMind}
\textbf{Use with EXAMPLES 3 \& 4}
Q: Model with Mathematics If the graph of a parent function is vertically stretched by a factor of 2 and then translated 3 units down, would you get the same graph if you translated the parent graph 3 units down first and then vertically stretched by a factor of 2? Explain.
\answer{No; the order the transformations are performed in affects the final result.}
\te{Examples are available at SavvasRealize.com.}
\end{te-addendum}

% Source page 31: 18 STEP 2 Understand & Apply
\begin{te-addendum}{5}{GeneralizeETP}
\textbf{EXAMPLE 5 Identify Transformations From an Equation}
\textbf{Use and Connect Mathematical Representations ETP—PART A}
Q: How are vertical transformations expressed differently in function notation than horizontal transformations?
\answer{When (f(x)) is written as (af(x)+k), (a) signals a vertical stretch or compression and (k) signals a vertical shift. When (f(x)) is written as (f[b(x-h)]), (b) signals a horizontal stretch or compression and (h) signals a horizontal shift.}
\textbf{PART B}
Q: What can you tell about the shape of the graph of (g) compared to the shape of the graph of (f), just by looking at the equation for (g)? Explain.
\answer{The shape of the graphs will be the same because the equation for (g) shows only a horizontal and a vertical translation, with no stretch or compression.}
\end{te-addendum}
\begin{example}{5}
\textbf{Identify Transformations From an Equation}
What transformations of (f(x)=x^2) result in the graph of the function (g)?
A. (g(x)=-(\frac{1}{3}x)^2)
(h(x)=(\frac{1}{3}x)^2=f(\frac{1}{3}x)) represents a horizontal stretch of the graph of (f) by the factor (\frac{1}{3}).
(g(x)=-(\frac{1}{3}x)^2=-h(x)=-f(\frac{1}{3}x)) represents a reflection across the x-axis of (f(\frac{1}{3}x)).
The graph of (g) is a horizontal stretch by the factor (\frac{1}{3}) and a reflection across the x-axis of the graph of (f).
\placeholder{graph}{Example 5A: blue upward parabola f=x^2, green dashed upward parabola h=x^2/9, red downward parabola g=-x^2/9, all vertex (0,0). All ends arrowed; color labels f,h,g. x grid spacing 1 with labels -6,-4,-2,2,4,6; y grid spacing 2 with labels -12,-8,-4,4,8,12. x/y axes arrowed and origin O; window approximately x=-7 to 7,y=-14 to 14. No point circles.}
\answer{A. A horizontal stretch by the factor (\frac{1}{3}) and a reflection across the x-axis.}
\te{Printed horizontal stretch factor 1/3; likely intended horizontal distances multiplied by 3. The printed coefficient-based wording is preserved.}
B. (g(x)=(x-4)^2+5)
(h(x)=(x-4)^2=f(x-4)) represents a translation 4 units to the right of the graph of (f).
(g(x)=(x-4)^2+5=h(x)+5=f(x-4)+5) represents a translation 5 units up of the graph of (f(x-4)).
The graph of (g) is a translation 4 units right and 5 units up of the graph of (f).
\placeholder{graph}{Example 5B: blue f=x^2 vertex (0,0), green dashed h=(x-4)^2 vertex (4,0), red g=(x-4)^2+5 vertex (4,5). Upward arrows on all visible curve ends, color labels f,h,g. x/y axes arrowed and origin O; grid spacing 1 in x,2 in y; x labels -6,-4,-2,2,4,6, y labels -12,-8,-4,4,8,12; window approximately x=-7 to 7,y=-14 to 14. No point circles.}
\answer{B. Translation 4 units right and 5 units up.}
\te{USE APPROPRIATE TOOLS: You can use graphing technology to graph the original and transformed equations to check that the transformations you have identified are correct.}
\end{example}
\begin{tryit}{5}
What transformations of the graph of (f(x)=|x|) are applied to graph the function (g)?
a. (g(x)=\frac{1}{2}|x+3|)
b. (g(x)=-|x|+2)
\answer{a. translation left 3 units and then vertical compression by the factor (\frac{1}{2}). b. reflection across the x-axis and then translation up 2 units.}
\end{tryit}
\begin{te-addendum}{5}{ElicitEvidence}
\textbf{Elicit and Use Evidence of Student Thinking ETP}
Q: What determines whether the function is transformed vertically or horizontally?
\answer{Vertical transformations are indicated by operations performed on the output of the function; horizontal transformations are indicated by operations performed on the input of the function.}
\te{Examples are available at SavvasRealize.com.}
\end{te-addendum}

% Source page 32: 19 STEP 2 Understand & Apply
\begin{te-addendum}{6}{PurposefulQuestions}
\textbf{EXAMPLE 6 Write an Equation From a Graph}
\textbf{Pose Purposeful Questions ETP}
Q: What do you notice about the graph of this function when compared to the parent function?
\answer{The graph of the parent function has been reflected vertically, shifted to the right, and shifted up.}
Q: How can you verify the value of (a) from the graph?
\answer{If the value of (a) is 1, the graph of the parent function can be placed on top of the transformed function to match its shape.}
\end{te-addendum}
\begin{example}{6}
\textbf{Write an Equation From a Graph}
\textbf{APPLICATION}
A scenic train ride makes trips on an old mining line. The graph shows the distance (y) in kilometers of the train from the station (x) minutes after the ride begins. What equation represents the distance from the station as a function of time? What is its domain?
\placeholder{graph}{Train graph: orange inverted V joins (0,0),(15,15),(30,0), vertex labeled (15,15). No point circles or curve arrows. Horizontal x-axis Time Since Start (min), labels 6,12,18,24,30,36, grid spacing 3. Vertical y-axis Distance From Station (km), labels 4,8,12,16,20, grid spacing 2. Axes arrowed, origin O, window x=0 to 36,y=0 to 20.}
\placeholder{photo}{Black steam locomotive with smoke against blue sky, viewed from the front-right, beside the train distance graph.}
The graph shows a reflection of an absolute value graph across the x-axis and a translation upward and to the right. The general form of this absolute value function is (y=-a|x-h|+k), where the point ((h,k)) represents the vertex and (-a) indicates that the graph opens downward.
Substituting the point of the vertex ((15,15)) for ((h,k)) gives the equation (y=-a|x-15|+15).
To solve for (a), you can use any point on the graph. Using the point ((0,0)) to substitute for ((x,y)) in the equation simplifies the computation:
(y=-a|x-15|+15)
(0=-a|0-15|+15)
(0=-15a+15)
(-15=-15a)
(a=1)
Now you can write the equation for distance as a function of time:
(y=-|x-15|+15).
According to the graph, the train returns to its station after 30 minutes, so the function’s domain is ([0,30]).
\answer{(y=-|x-15|+15); domain ([0,30]).}
\te{STUDY TIP: Solving algebraically is only one method for determining a stretch or compression factor.}
\end{example}
\begin{tryit}{6}
How would the graph and equation be affected if the train traveled twice as far in the same amount of time?
\answer{There would be a vertical stretch by a factor of 2. The new equation would be (y=2(-|x-15|+15)=-2|x-15|+30).}
\end{tryit}
\begin{te-addendum}{6}{ElicitEvidence}
\textbf{Elicit and Use Evidence of Student Thinking ETP}
Q: How can the phrase “traveled twice as far in the same amount of time” be interpreted with regard to the domain and range?
\answer{The domain will stay the same because there was no change in time, and the range will be doubled.}
\end{te-addendum}
\begin{te-addendum}{6}{HabitsOfMind}
\textbf{Use with EXAMPLES 5 \& 6}
Q: Make Sense and Persevere The function (f(x)=|x|) is translated 3 units right and 2 units down, and then vertically stretched by a factor of 4. What is an equation for the transformed function (g)?
\answer{(g(x)=4(|x-3|-2)), or (g(x)=4|x-3|-8)}
\end{te-addendum}
\begin{te-addendum}{6}{RtIExtend}
\textbf{Extend Student Thinking}
\textbf{USE WITH EXAMPLE 6} Students use what they know about writing an equation from a graph to practice writing an equation from context.
\item Another train ride is available that lasts 45 minutes and travels a total distance of 25 km.
Q: Explain how to write an equation to model this situation.
\answer{Because the total time is 45 min and the total distance is 25 km, the turn-around point, or vertex, will occur at ((22.5,12.5)). Plug the values of the vertex into the equation for (h) and (k), along with the point ((0,0)), and solve for (a).}
Q: What equation represents the distance from the station as a function of time?
\answer{(d=-\frac{5}{9}|m-22.5|+12.5)}
\te{Examples are available at SavvasRealize.com.}
\end{te-addendum}


% Source page 33: 20 STEP 2 Understand & Apply
\begin{te-addendum}{ConceptSummary}{PurposefulQuestions}
\textbf{CONCEPT SUMMARY Transformations of Functions}
Q: What happens to the graph of (f(x)=a\cdot f[b(x-h)]+k) as you change the values of (a), (b), (h), and (k)?
\answer{Changing (a) results in a vertical stretch or compression. Changing (b) results in a horizontal stretch or compression. Changing (h) results in a horizontal translation. Changing (k) results in a vertical translation.}
\end{te-addendum}
\begin{te-addendum}{ConceptSummary}{CommonError}
\textbf{Do You UNDERSTAND? | Do You KNOW HOW?}
\textbf{Common Error—Exercise 5} Students may incorrectly interpret the negative sign in the function (g) as a translation of 3 units to the left. When students are unsure, suggest that they make a table of values to check the points on their graph.
\end{te-addendum}
\begin{concept-summary}
\textbf{CONCEPT SUMMARY Transformations of Functions}
For a function (f(x)), the graph of (f(x)=a\cdot f[b(x-h)]+k) represents a transformation of the graph of that function by translation, reflection, or stretching.
\begin{tabular}{ll}
WORDS & EQUATIONS \\
Horizontal translation of (f) right 2 units (altering (h)) & (g(x)=f(x-2)); (f(x)=x^2+x); (g(x)=(x-2)^2+(x-2)=x^2-3x+2) \\
Vertical translation of (f) up 3 units (altering (k)) & (h(x)=f(x)+3); (f(x)=x^2+x); (h(x)=x^2+x+3) \\
Reflection of (f) across the x-axis (altering (a)) & (g(x)=-f(x)); (f(x)=x^2+x); (-f(x)=-(x^2+x)=-x^2-x) \\
Reflection of (f) across the y-axis (altering (b)) & (h(x)=f(-x)); (f(x)=x^2+x); (f(-x)=(-x)^2+(-x)=x^2-x) \\
Horizontal stretch of (f) by the factor (\frac{1}{2}) (altering (b)) & (g(x)=f(\frac{1}{2}x)); (f(x)=x^2+x); (f(\frac{1}{2}x)=(\frac{1}{2}x)^2+\frac{1}{2}x=\frac{1}{4}x^2+\frac{1}{2}x) \\
Vertical stretch of (f) by the factor 2 (altering (a)) & (h(x)=2f(x)); (f(x)=x^2+x); (2f(x)=2(x^2+x)=2x^2+2x) \\
\end{tabular}
\te{Printed horizontal stretch factor 1/2; likely describes the input coefficient, with horizontal distances multiplied by 2.}
\placeholder{graph}{Summary translation graph: red parabola f=x^2+x, blue g=x^2-3x+2, green h=x^2+x+3. Vertices (-0.5,-0.25),(1.5,-0.25),(-0.5,2.75). All upper ends arrowed, no point circles or curve labels. Unit grid, x/y axes arrowed, origin O, x labels -4,-2,2,4, y labels -2,2,4,6; window approximately x=-5 to 5,y=-3 to 7.}
\placeholder{graph}{Summary reflection graph: red upward f=x^2+x, blue downward -x^2-x, green upward x^2-x. Vertices (-0.5,-0.25),(-0.5,0.25),(0.5,-0.25). All curve ends arrowed. Unit grid, x/y axes arrowed and origin O; labels -4,-2,2,4 on each axis; window approximately [-5,5] by [-5,5]. No point circles.}
\placeholder{graph}{Summary stretch graph: red f=x^2+x, blue x^2/4+x/2, green 2x^2+2x, upward parabolas with vertices (-0.5,-0.25),(-1,-0.25),(-0.5,-0.5). All upper ends arrowed. Unit grid, x/y axes arrowed, origin O, labels -4,-2,2,4 on each axis, window approximately [-5,5] by [-5,5]. No point circles.}
\textbf{Do You UNDERSTAND?}
1. ESSENTIAL QUESTION How do the values (a), (b), (h), and (k) in (a\cdot f[b(x-h)]+k) transform the graph of a parent function, (f(x))?
\answer{From the equations, you can see whether the graph of the parent function will be shifted up or down, reflected over an axis, or stretched or compressed to create the graph of the new function.}
2. Reason Do (k) and (h) affect the input or output for (g(x)=f(x)+k) and (g(x)=f(x-h))? Explain.
\answer{(g(x)=f(x)+k) is a vertical translation (k) units, so the output of the function is affected. (g(x)=f(x-h)) is a horizontal translation (h) units, so the input of the function is affected.}
3. Error Analysis Margo is comparing the functions (f(x)=|x|) and (g(x)=|x+1|-5). She said the graph of (g) is a vertical translation of the graph of (f) 5 units down and a horizontal translation of the graph of (f) 1 unit right. What is Margo’s error?
\answer{(g(x)) is a horizontal translation of (f(x)) left 1 unit, not right 1 unit, because (h<0).}
\textbf{Do You KNOW HOW?}
Graph each function and its parent function.
4. (g(x)=|x|-1)
\answer{See graph.}
\placeholder{graph}{Answer 4: blue f=|x| vertex (0,0), red g=|x|-1 vertex (0,-1), roots -1,1. Unit square grid [-5,5] by [-5,5]; x/y axes arrowed, origin O, labels -4,-2,2,4 on each axis; all curve ends arrowed, color labels f,g, no point circles.}
5. (g(x)=(x-3)^2)
\answer{See graph.}
\placeholder{graph}{Answer 5: blue f=x^2 vertex (0,0), red g=(x-3)^2 vertex (3,0). Unit square grid [-5,5] by [-5,5], arrowed x/y axes, origin O, labels -4,-2,2,4 on each axis. Upper curve ends arrowed, color labels f,g, no point circles.}
6. (g(x)=-|x|)
\answer{See graph.}
\placeholder{graph}{Answer 6: blue upward V f=|x| and red downward V g=-|x|, both vertex (0,0), slopes magnitude 1. Unit square grid [-5,5] by [-5,5], arrowed x/y axes, origin O, labels -4,-2,2,4 on both axes. Curve ends arrowed, color labels f,g, no point circles.}
7. (g(x)=-x)
\answer{See graph.}
\placeholder{graph}{Answer 7: blue line f=x and red line g=-x through origin, arrows on both ends. Unit square grid [-5,5] by [-5,5], arrowed x/y axes, origin O, labels -4,-2,2,4 on both axes. Color labels f,g, no point circles.}
8. (g(x)=x^2-2)
\answer{See graph.}
\placeholder{graph}{Answer 8: blue parabola f=x^2 vertex (0,0), red g=x^2-2 vertex (0,-2). Unit square grid [-5,5] by [-5,5], arrowed x/y axes, origin O, labels -4,-2,2,4 on both axes. Upper curve ends arrowed, color labels f,g, no point circles.}
9. (g(x)=\frac{1}{2}|x|)
\answer{See graph.}
\placeholder{graph}{Answer 9: blue V f=|x| and red shallower V g=|x|/2 with common vertex (0,0). Unit square grid [-5,5] by [-5,5], arrowed x/y axes, origin O, labels -4,-2,2,4 on both axes. Upper curve ends arrowed, color labels f,g, no point circles.}
10. (g(x)=4x)
\answer{See graph.}
\placeholder{graph}{Answer 10: blue line f=x and red steep line g=4x, both through origin. Unit square grid [-5,5] by [-5,5], arrowed x/y axes, origin O, labels -4,-2,2,4 on both axes. Lines arrowed at both ends, color labels f,g, no point circles.}
11. (g(x)=|5x|)
\answer{See graph.}
\placeholder{graph}{Answer 11: blue V f=|x| and steep red V g=|5x|, both vertex (0,0). Unit square grid [-5,5] by [-5,5], arrowed x/y axes, origin O, labels -4,-2,2,4 on both axes. Curve ends arrowed, color labels f,g, no point circles.}
12. Which types of transformations in Exercises 4--11 do not change the shape of a graph? Which types of transformations change the shape of a graph? Explain.
\answer{Vertical and horizontal translations and reflections do not change the shape of a graph. Stretches and compressions change the shape of a graph because they affect the steepness of the graph.}
\te{Concept Summary, Do You Understand, and Do You Know How are available at SavvasRealize.com.}
\end{concept-summary}


% Source page 34: 21 STEP 3 Practice & Problem Solving
\begin{te-addendum}{Practice}{GeneralizeETP}
\textbf{PRACTICE \& PROBLEM SOLVING}
\textbf{Lesson Practice} You may opt to have students complete the automatically scored Practice and Problem Solving items online powered by MathXL for School.
Choose from: Lesson Practice; Adaptive Practice.
You may also take advantage of the bank of exercises for assigning additional practice.
\textbf{Assignment Guide}
\begin{tabular}{ll}
On-Level & Advanced \\
14, 16--38 & 13--38 \\
\end{tabular}
\textbf{Engage Through Student Choice}
Promote student agency by allowing students to choose practice items. You may structure this choice in many ways.
For example:
Assign each section a point value. Students choose at least one item from each section and items chosen should have a minimum of 20 total points.
\begin{tabular}{ll}
Understand, Apply & 2 points each \\
Practice & 1 point each \\
Assessment Practice & 1 point each \\
Performance Task & 3 points \\
\end{tabular}
\te{Practice \& Problem Solving and video tutorials are available at SavvasRealize.com. Scan the QR code to access a video tutorial.}
\end{te-addendum}
\begin{item-analysis}
\begin{tabular}{lll}
\toprule
Example & Items & DOK \\
\midrule
1 & 18--21 & 1 \\
1 & 13 & 3 \\
2 & 16, 22, 23 & 1 \\
2 & 14 & 2 \\
3 & 24--27 & 1 \\
3 & 15 & 3 \\
4 & 28 & 1 \\
5 & 29--32, 36 & 2 \\
5 & 17, 37 & 3 \\
6 & 33--35, 38 & 3 \\
\bottomrule
\end{tabular}
\end{item-analysis}
\begin{practice}{13}{3}
UNDERSTAND. Use Structure Write a function (g) with the parent function (f(x)=x^2) that has a vertex at ((3,-6)).
\answer{(g(x)=(x-3)^2-6)}
\end{practice}
\begin{practice}{14}{2}
Error Analysis Describe and correct the error a student made in graphing (g(x)=f(-x)) as a reflection across the y-axis of the graph of (f(x)=|x+2|+1).
\placeholder{graph}{Student error graph: blue upward V f with vertex (-2,1), slopes magnitude 1, red downward V g with vertex (-2,-1), slopes magnitude 1. All ends arrowed; color labels f,g, large red X at lower right. Unit pale-blue grid, gray arrowed x/y axes, origin O, x labels -4,-2,2,4, y labels -4,2,4; window approximately [-6,6] by [-6,6].}
\answer{The student reflected the graph of (f(x)) across the x-axis instead of reflecting the graph across the y-axis.}
\placeholder{graph}{Printed corrected answer 14: red upward V f=|x+2|+1, vertex (-2,1); blue upward V g=|x-2|+1, vertex (2,1); both cross y-axis at 3. Unit square grid [-5,5] by [-5,5], arrowed x/y axes, origin O, labels -4,-2,2,4 on each axis; all curve ends arrowed, no point circles, color labels f,g.}
\end{practice}
\begin{practice}{15}{3}
Higher Order Thinking Describe the transformation (g) of (f(x)=|x|) as a stretch and as a compression. Then write two equations to represent the function. What can you conclude? Explain.
\placeholder{graph}{Red upward V g=4|x|, vertex (0,0), arrows on upper ends. x/y axes arrowed, origin O, x unit grid and labels -4,-2,2,4; y grid spacing 2 and labels 4,8,12,16; window approximately x=-5 to 5,y=-4 to 20. No point circles.}
\answer{Multiplying the y-values of the points on the graph of (f) by 4 results in points on the graph of (g). So, (g(x)) is a vertical stretch of (f(x)) by a scale factor of 4 and can be written as (g(x)=4|x|). Dividing the x-values of the points on the graph of (f) also results in points on the graph of (g). So (g(x)) is a horizontal compression of (f(x)) by a scale factor of 4 and can be written as (g(x)=|4x|). A vertical stretch is the same as a horizontal compression in this case.}
\te{Printed “Dividing the x-values” omits a divisor; likely “by 4.”}
\end{practice}
\begin{practice}{16}{1}
Use Structure The graph of the parent function (f(x)=x^2) is reflected across the y-axis. Write an equation for the function (g) after the reflection. Show your work. Based on your equation, what happens to the graph? Explain.
\answer{(g(x)=f(-x)=(-x)^2=(x)^2); The graph of (g(x)) is the same as the graph of (f(x)).}
\te{Answer printed on PDF page 35.}
\end{practice}
\begin{practice}{17}{3}
Error Analysis Monisha is comparing (f(x)=|x|) and (g(x)=|2x-4|). She said the graph of (g) is a horizontal translation of the graph of (f) 4 units to the right and a horizontal compression of the graph of (f) by a factor of 2. What is Monisha’s error?
\answer{(g(x)=|2x-4|) is equivalent to (g(x)=2|x-2|). So, (g(x)) is a horizontal translation of (f(x)) right 2 units and a vertical stretch of (f(x)) by a factor of 2.}
\te{Answer printed on PDF page 35.}
\end{practice}
\begin{practice}{18}{1}
PRACTICE. Graph each function as a translation of its parent function, (f). How did the transformation affect the domain and range? SEE EXAMPLE 1
(g(x)=|x|-5)
\answer{The domain values of (f) and (g) are the same, but the range values of (g) are 5 units less than the range values of (f).}
\placeholder{graph}{Answer 18 on PDF35: blue V f=|x| vertex (0,0), red V g=|x|-5 vertex (0,-5), roots -5,5. Unit grid [-5,5] by [-5,5], x/y axes arrowed, origin O, labels -4,-2,2,4 on both axes; curve ends arrowed and color labels f,g; no point circles.}
\end{practice}
\begin{practice}{19}{1}
Graph each function as a translation of its parent function, (f). How did the transformation affect the domain and range? SEE EXAMPLE 1
(g(x)=(x+1)^2)
\answer{The range values of (f) and (g) are the same, but the domain values of (g) are 1 unit less than the domain values of (f) at corresponding range values.}
\placeholder{graph}{Answer 19 on PDF35: blue upward parabola f=x^2 vertex (0,0), red g=(x+1)^2 vertex (-1,0). Unit grid [-5,5] by [-5,5], x/y axes arrowed, origin O, labels -4,-2,2,4 on both axes; upper curve ends arrowed, color labels f,g, no point circles.}
\end{practice}
\begin{practice}{20}{1}
Graph each function as a translation of its parent function, (f). How did the transformation affect the domain and range? SEE EXAMPLE 1
(g(x)=|x-3|)
\answer{The range values of (f) and (g) are the same, but the domain values of (g) are 3 units greater than the domain values of (f) at corresponding range values.}
\placeholder{graph}{Answer 20 on PDF35: blue V f=|x| vertex (0,0), red V g=|x-3| vertex (3,0), y-intercept 3. Unit grid [-5,5] by [-5,5], x/y axes arrowed, origin O, labels -4,-2,2,4 on both axes; all curve ends arrowed, color labels f,g, no point circles.}
\end{practice}
\begin{practice}{21}{1}
Graph each function as a translation of its parent function, (f). How did the transformation affect the domain and range? SEE EXAMPLE 1
(g(x)=x^2+2)
\answer{The domain values of (f) and (g) are the same, but the range values of (g) are 2 units greater than those of (f).}
\placeholder{graph}{Answer 21 on PDF35: blue upward parabola f=x^2 vertex (0,0), red g=x^2+2 vertex (0,2). Unit grid [-5,5] by [-5,5], x/y axes arrowed, origin O, labels -4,-2,2,4 on both axes; all upper curve ends arrowed, color labels f,g, no point circles.}
\end{practice}
\begin{practice}{22}{1}
What is the function for the image graph? Check by graphing. SEE EXAMPLE 2
Reflect (f(x)=x^2+1) across the x-axis.
\answer{(f(x)=-x^2-1)}
\end{practice}
\begin{practice}{23}{1}
What is the function for the image graph? Check by graphing. SEE EXAMPLE 2
Reflect (f(x)=x^2+1) across the y-axis.
\answer{(f(x)=x^2+1)}
\end{practice}
\begin{practice}{24}{1}
Graph each function as a vertical stretch or compression of its parent function. SEE EXAMPLE 3
(g(x)=0.25|x|)
\answer{See back of book.}
\end{practice}
\begin{practice}{25}{1}
Graph each function as a vertical stretch or compression of its parent function. SEE EXAMPLE 3
(g(x)=3x^2)
\answer{See back of book.}
\end{practice}
\begin{practice}{26}{1}
Graph each function as a vertical stretch or compression of its parent function. SEE EXAMPLE 3
(g(x)=1.5|x|)
\answer{See back of book.}
\end{practice}
\begin{practice}{27}{1}
Graph each function as a vertical stretch or compression of its parent function. SEE EXAMPLE 3
(g(x)=0.75x^2)
\answer{See back of book.}
\end{practice}
\begin{practice}{28}{1}
Use the graph of (f(x)) to graph (y=f(x+1)+2). SEE EXAMPLE 4
\placeholder{graph}{Blue polyline f joins filled endpoint (-3,0), vertex (1,-4), vertex (2,3), filled endpoint (5,0). Unit square grid [-5,5] by [-5,5]; x/y axes arrowed, origin O, x labels -4,2,4 and y labels -4,2,4. No curve arrows.}
\answer{See back of book.}
\end{practice}
\begin{practice}{29}{2}
What transformations of (f(x)=x^2) are applied to the function (g)? Check by graphing. SEE EXAMPLE 5
(g(x)=2(x+1)^2)
\answer{See back of book.}
\end{practice}
\begin{practice}{30}{2}
What transformations of (f(x)=x^2) are applied to the function (g)? Check by graphing. SEE EXAMPLE 5
(g(x)=(x-3)^2+5)
\answer{See back of book.}
\end{practice}
\begin{practice}{31}{2}
What transformations of (f(x)=x^2) are applied to the function (g)? Check by graphing. SEE EXAMPLE 5
(g(x)=-x^2-6)
\answer{See back of book.}
\end{practice}
\begin{practice}{32}{2}
What transformations of (f(x)=x^2) are applied to the function (g)? Check by graphing. SEE EXAMPLE 5
(g(x)=4(x-7)^2-9)
\answer{See back of book.}
\end{practice}
\begin{practice}{33}{3}
The graph shows the height (y) in feet of a flying insect (x) seconds after taking off from the ground. Write an equation that represents the height of the insect as a function of time. SEE EXAMPLE 6
\placeholder{graph}{Red downward parabola y=-(x-3)^2+9, visible from (0,0) to (6,0), vertex (3,9), downward-right arrow near (6,0), no marked point circles. x-axis Time (s), labels 0,2,4,6, unit grid; y-axis Height (ft), labels 0,2,4,6,8, unit grid. Axes arrowed, window approximately [0,7] by [0,9].}
\answer{(y=-(x-3)^2+9)}
\end{practice}

% Source page 35: 22 STEP 3 Practice & Problem Solving
\begin{te-addendum}{Practice}{GeneralizeETP}
\textbf{Answers}
\te{Printed answers and graphs for Exercises 16–21 are attached to their corresponding practice blocks above. The repeated line “24–32, 34. See back of book.” is represented in those answer blocks. Practice \& Problem Solving is available at SavvasRealize.com.}
\end{te-addendum}
\begin{practice}{34}{3}
APPLY. Model With Mathematics Chiang walks to school each day. She passes the library halfway on her walk to school. She walks at a rate of 1 block per minute. The graph shows the distance Chiang is from the library as she walks to school.
\placeholder{graph}{Blue V y=|x| vertex (0,0), arrowheads upper left/right, unit grid approximately x=-5 to 5,y=-1 to 8. Horizontal x-axis labels -4,-2,2,4, left caption Minutes Before Passing Library, right caption Minutes After Passing Library. Vertical y-axis Blocks From Library, labels 2,4,6. Axes arrowed, origin O, no point circles.}
a. Write a function, (f), to model the distance Chiang is from the library when she walks to school.
b. If Chiang jogs to school, she travels at a rate of 2.5 blocks per minute. Write a function, (g), to model the distance Chiang is from the library when she jogs to school.
c. Graph the function, (g), that models the distance Chiang is from the library when she jogs to school.
\answer{See back of book.}
\end{practice}
\begin{practice}{35}{3}
Model With Mathematics The archer fish spits stream of water at flying insects to knock them into the lake. The path of the stream of water is shown with (x) and (y) distances in feet. Write an equation to represent the path of the stream of water in relation to the coordinate grid. Then determine the coordinates of the point of maximum height of the stream of water.
\placeholder{graph}{Black downward parabola through filled (0,0),(1,1),(2,0), vertex (1,1), both descending ends arrowed, over illustration of fish spitting at insect on leaf. x/y axes arrowed, origin 0, x labels 1,2, y labels -1,1; half-unit grid, window approximately x=-1 to 3,y=-1.5 to 1.5.}
\answer{(y=-(x-1)^2+1); (1,1)}
\end{practice}
\begin{practice}{36}{2}
ASSESSMENT PRACTICE. Match each equation on the left to an equation on the right that has the same translation of (y=|x|).
A. (y=|x|-1)
B. (y=-|x+1|)
C. (y=-|x+1|-1)
D. (y=|x+1|-1)
E. (y=2|x|-1)
F. (y=|x+1|)
\answer{A–E; B–F; C–D.}
\te{Matches are printed as magenta connecting lines.}
\end{practice}
\begin{practice}{37}{3}
SAT/ACT Which translation is part of transforming (f(x)=x^2) into (h(x)=(x+4)^2-2)?
A. left 4 units
B. left 2 units
C. right 2 units
D. right 4 units
\answer{A}
\end{practice}
\begin{practice}{38}{3}
Performance Task The Louvre Pyramid in Paris is shown on the coordinate grid, where (x) and (y) are measured in meters and the ground is represented by the x-axis.
\placeholder{graph}{Illustration of Louvre Pyramid, black outline joins (0,0), filled apex (17.7,21.6), filled right endpoint (35.4,0), with callouts (17.7,21.6) and (35.4,0). x/y axes arrowed, x labels 0,20,40 and y label 20, square-looking grid spacing 5 meters. Visible window approximately x=0 to 40,y=0 to 25. No arrows on pyramid outline.}
Part A The outline of the Pyramid is a transformation of the function (f(x)=|x|). Write a function (g) to model the outline of the Pyramid.
\answer{Part A (g(x)=-1.22|x-17.7|+21.6)}
Part B What is the domain and range of the function that models the outline of the Pyramid? What do the domain and range represent?
\answer{Part B domain: (0\leq x\leq35.4); range: (0\leq y\leq21.6); The domain represents the width of the pyramid, which is 35.4 m wide. The range represents the height of the pyramid, which is 21.6 m high.}
\end{practice}


% Source page 36: 22A STEP 4 Assess & Differentiate
\begin{te-addendum}{Assess}{RtISupport}
\textbf{LESSON QUIZ}
Use the Lesson Quiz to assess students’ understanding of the mathematics in the lesson.
Students can take the Lesson Quiz online or you can download a printable copy from SavvasRealize.com. The Lesson Quiz is also available in the Assessment Sourcebook.
\textbf{Item Analysis}
\begin{tabular}{lll}
Item & DOK & Skills Review and Practice \\
1 & 2 & F68 \\
2 & 1 & F57 \\
3 & 1 & F57 \\
4 & 2 & F56 \\
5 & 2 & F25, F57 \\
\end{tabular}
Use the student scores on the Lesson Quiz to prescribe differentiated assignments.
If students take the Lesson Quiz online, it will be automatically scored and appropriate differentiated practice will be assigned based on student performance.
\begin{tabular}{lll}
Intervention & 0--3 points & Reteach to Build Understanding; Mathematical Literacy and Vocabulary; Lesson Virtual Nerd Videos \\
On-Level & 4 points & Enrichment \\
Advanced & 5 points & Enrichment \\
\end{tabular}
\end{te-addendum}
\begin{te-addendum}{Assess}{GeneralizeETP}
\textbf{ASSESSMENT RESOURCES}
\textbf{1-2 Lesson Quiz: Transformations of Functions}
Name blank.
1. Select all the transformations of (f(x)=x^2) that combine to result in the graph of function (g) at the right.
\placeholder{graph}{Quiz 1: black downward parabola with vertex (-1,-2), consistent with g=-(x+1)^2/4-2. No x-intercepts or marked points; arrows on both descending ends. Unit square grid approximately x=-5 to 5,y=-7 to 3; arrowed x/y axes, origin O; x labels -4,-2,2,4, y labels -6,-4,-2,2.}
A. translation of 1 unit left
B. translation of 1 unit right
C. translation of 2 units down
D. horizontal stretch by a factor of 2
E. vertical stretch by a factor of 2
F. reflection across the x-axis
\answer{1. A, C, D, F.}
2. What is the equation for (g), which is (f(x)=2x^2+3x-1) reflected across the y-axis?
A. (g(x)=2x^2+3x-1)
B. (g(x)=-2x^2-3x+1)
C. (g(x)=2x^2-3x-1)
D. (g(x)=-2x^2-3x-1)
\answer{2. C}
3. The graph represents (f). Graph (g(x)=2f(x-1)).
\placeholder{graph}{Quiz 3: black f polyline (-2,2),(-1,-1),(1,-1),(2,0). Printed magenta answer g joins (-1,4),(0,-2),(2,-2),(3,0), labeled g; black f label near left segment. Unit square grid [-5,5] by [-5,5], arrowed x/y axes, origin O, labels -4,-2,2,4 on both axes; no point circles or curve arrows.}
\answer{3. See graph.}
4. A bucket has a leak that started 6 hours ago. At time 0, when the bucket just emptied, a tap was turned on for 3 hours. The amount of water in the bucket can be modeled by the function (f(x)=|x|) over the interval ([-6,3]), where (x) is the time, in hours, and (y) is the amount of water in the bucket, in gallons. A similarly full bucket started to leak only 3 hours ago. Which translation of the graph represents the amount of water in the second bucket if it is moved under the tap as soon as it empties?
A. (g(x)=|x|+3)
B. (g(x)=|x+3|)
C. (g(x)=|x|-3)
D. (g(x)=|x-3|)
\answer{4. D}
5. The graph of the function (f) is shown. The domain of (f) is ([0,50]). What is the range of (g(x)=4f(x))?
\placeholder{graph}{Quiz 5: black inverted V joins (0,0),(25,1.75),(50,0), no point circles or line arrows. x-axis labels 10,20,30,40,50, grid spacing 5; y-axis labels 0.5,1,1.5, grid spacing 0.25. Arrowed x/y axes, origin O, window approximately x=-5 to 60,y=-0.25 to 2.}
A. ([0,50])
B. ([0,70])
C. ([0,1.75])
D. ([0,7])
\answer{5. D}
\te{Lesson Quiz is available at SavvasRealize.com. enVision Algebra 2 • Assessment Sourcebook. Copyright © Savvas Learning Company LLC. All Rights Reserved.}
\end{te-addendum}

% Source page 37: 22B STEP 4 Assess & Differentiate
\begin{te-addendum}{Assess}{GeneralizeETP}
\textbf{DIFFERENTIATED RESOURCES}
I = Intervention; O = On-Level; A = Advanced.
\textbf{Reteach to Build Understanding—I}
Provides scaffolded reteaching for the key lesson concepts.
MathXL for School.
\textbf{1-2 Reteach to Build Understanding: Transformations of Functions}
Name blank.
If you are given a parent function, such as (f), and positive constants (a), (k), and (h), you can transform the function using the following transformations.
\begin{tabular}{ll}
(g(x)=f(x)+k) & Vertical translation up; shifts the graph of (f) up (k) units. \\
(g(x)=f(x)-k) & Vertical translation down; shifts the graph of (f) down (k) units. \\
(g(x)=f(x+h)) & Horizontal translation left; shifts the graph of (f) left (h) units. \\
(g(x)=f(x-h)) & Horizontal translation right; shifts the graph of (f) right (h) units. \\
(g(x)=f(-x)) & Reflection across the y-axis; reverses the sign of the x-coordinate. \\
(g(x)=-f(x)) & Reflection across the x-axis; reverses the sign of the y-coordinate. \\
(g(x)=f(ax)) & Horizontal stretch or compression of (f) by the factor (a); the graph moves closer to or farther from the y-axis. \\
(g(x)=af(x)) & Vertical stretch or compression of (f) by the factor (a); the graph moves closer to or farther from the x-axis. \\
\end{tabular}
1. Match the following equations with the transformation that occurred with a parent function of (f).
a. (g(x)=f(x)+4)
b. (g(x)=f(x)-4)
c. (g(x)=f(x+4))
d. (g(x)=f(x-4))
e. (g(x)=f(-x))
f. (g(x)=-f(x))
g. (g(x)=f(4x))
h. (g(x)=4f(x))
Reflection across the y-axis; Horizontal translation left; Vertical stretch; Vertical translation down; Vertical translation up; Horizontal compression; Horizontal translation right; Reflection across the x-axis.
\answer{a. Vertical translation up; b. Vertical translation down; c. Horizontal translation left; d. Horizontal translation right; e. Reflection across the y-axis; f. Reflection across the x-axis; g. Horizontal compression; h. Vertical stretch.}
\te{Answers are printed as magenta connecting lines.}
2. Beth said that (g(x)=f(x)=x^2-12) is a horizontal translation of (f(x)=x^2). Find and fix the errors and write the correct equation for a horizontal translation.
\answer{Sample answer: Beth should have said it was a vertical translation down. If it were a horizontal translation, it should be (f(x)=(x-12)^2).}
3. Use the parent function (f(x)=x^2) to fill in the blanks.
a. (f(x)=-x^2) is a blank across the x-axis.
\answer{reflection}
b. (f(x)=\text{blank}) is a horizontal translation right by 3 units.
\answer{(x-3)^2}
c. (f(x)=x^2+6) is a vertical translation up by blank.
\answer{6 units}
\end{te-addendum}
\begin{te-addendum}{Assess}{GeneralizeETP}
\textbf{Additional Practice—I, O}
Provides extra practice for each lesson.
MathXL for School.
\textbf{1-2 Additional Practice: Transformations of Functions}
Name blank.
1. Graph the function (g(x)=|x|+4) as a translation of the parent function (f) shown. How did the transformation affect the domain and range?
\placeholder{graph}{Additional Practice 1: black V f(x)=|x| vertex (0,0), magenta g(x)=|x|+4 vertex (0,4); upper curve ends arrowed. Unit grid approximately x=-5 to 5,y=0 to 9; x/y axes arrowed, origin O, x labels -4,-2,2,4, y labels 2,6,8. Curves labeled f(x),g(x); no point circles.}
\answer{Domain of (f(x)) and (g(x)) are the same. Range values are (f(x)=y\geq0) and (g(x)=y\geq4).}
For Items 2 and 3, what is the equation for each reflected graph of (f(x)=x^2-4)?
2. Reflect across the x-axis.
\answer{(f(x)=-x^2+4)}
3. Reflect across the y-axis.
\answer{(f(x)=x^2-4)}
Graph each function as a vertical stretch or compression of the parent function (f).
4. (g(x)=4.5|x|)
\answer{See graph.}
\placeholder{graph}{Additional Practice 4 answer: black V f(x)=|x| and magenta steep V g(x)=4.5|x|, vertex (0,0), upper ends arrowed. Unit square grid approximately x=-5 to 5,y=-2 to 8; x/y axes arrowed, origin O, x labels -4,-2,2,4, y labels 4,6. Color labels f(x),g(x), no point circles.}
5. (g(x)=0.5x^2)
\answer{See graph.}
\placeholder{graph}{Additional Practice 5 answer: black upward parabola f(x)=x^2 and wider magenta g(x)=0.5x^2, common vertex (0,0). Unit grid approximately x=-5 to 5,y=-1 to 9, x/y axes arrowed, origin O, x labels -4,-2,2,4, y labels 2,4,6,8. Upper ends arrowed, color labels f(x),g(x), no point circles.}
What transformations of (f(x)=x^2) are applied to get the function (g)?
6. (g(x)=3(x+2)^2)
\answer{translation of 2 units to the left; vertical stretch by 3}
7. (g(x)=-(x-5)^2+1)
\answer{translation of 5 units to the right; reflection across x-axis; vertical translation up 1}
8. Derek walks to his best friend’s house at a rate of 1 block per minute, then turns around and walks home. The graph shows the distance Derek walks in the given amount of time. Write an equation for the graph.
\placeholder{graph}{Additional Practice 8: black inverted V from (0,0) to filled (10,10) to (20,0), endpoints appear filled. Horizontal axis Time (min), labels 0,4,8,12,16,20 with grid 2; vertical axis Distance From Home (blocks), labels 2,4,6,8,10, grid 1. Axes arrowed, window [0,20] by [0,10]; no curve arrows.}
\answer{(g(x)=-|x-10|+10)}
9. Given the parent function (f(x)=x^2), what is the new equation if the function is translated 4 units to the right and 3 units down?
\answer{(f(x)=(x-4)^2-3)}
\end{te-addendum}
\begin{te-addendum}{Assess}{GeneralizeETP}
\textbf{Enrichment—O, A}
Presents engaging problems and activities that extend the lesson concepts.
MathXL for School.
\textbf{1-2 Enrichment: Transformation of Functions}
Name blank.
For Item 1, write the function based on the translation of (f(x)=|x-1|). Then for Item 2, write the translation based on the function in Item 1. Continue making the translation based on the preceding item. Then, create a graph of the final function.
1. reflect across the x-axis
\answer{(f(x)=-|x-1|)}
2. reflect across the y-axis
\answer{(f(x)=-|-x-1|)}
3. compress vertically by a factor of (\frac{1}{5})
\answer{(f(x)=-\frac{1}{5}|-x-1|)}
4. translate 42 units down
\answer{(f(x)=-\frac{1}{5}|-x-1|-42)}
5. compress horizontally by a factor of (\frac{1}{12})
\answer{(f(x)=-\frac{1}{5}|-(12)x-1|-42)}
6. stretch vertically by a factor of 2
\answer{(f(x)=-\frac{2}{5}|-(12)x-1|-42)}
7. translate 23 units left
\answer{(f(x)=-\frac{2}{5}|-(12)x-22|-42)}
8. stretch horizontally by a factor of 7
\answer{(f(x)=-\frac{2}{5}|-(\frac{12}{7})x-22|-42)}
(f(x)=-\frac{2}{5}|-(\frac{12}{7})x-22|-42)
\placeholder{graph}{Enrichment final graph: magenta inverted V, vertex approximately (-12.8,-42), crosses y-axis approximately -51, tails to approximately (-70,-81) and (50,-85), with arrowheads. Grid spacing 10 on both axes, x labels -60,-40,-20,20,40, y labels 10,-20,-60,-80; arrowed x/y axes, origin O, visible window approximately x=-70 to 50,y=-90 to 20. No marked points.}
\te{Printed step 6 leaves the -42 unchanged under a vertical stretch; likely -84. Printed step 7 changes -1 to -22 for a 23-unit left translation; likely the input substitution should be applied to the entire preceding input expression. Printed equations and final graph are retained. The preview title prints singular “Transformation.”}
\end{te-addendum}
\begin{te-addendum}{Assess}{ELLAddendum}
\textbf{Mathematical Literacy and Vocabulary—I, O}
Helps students develop and reinforce understanding of key terms and concepts.
\textbf{1-2 Mathematical Literacy and Vocabulary: Transformations of Functions}
Name blank.
Fill in the missing words from the concept list to complete the web diagram.
\begin{tabular}{lll}
reflection & vertical translation & vertical compression \\
(af(x)), where (a>1) & (f(x+h)) & (-f(x)) \\
(af(x)), where (0<a<1) & & \\
\end{tabular}
\placeholder{diagram}{Web diagram: central hexagon “Transformations of the parent function f.” Five branches lead to rectangles: top reflection, left vertical translation, right horizontal translation, lower-left vertical stretch, lower-right vertical compression. Top rectangle branches to circles f(-x) and -f(x); left to circle f(x)+c; right to circle f(x+h); lower-left to circle af(x), where a>1; lower-right to circle af(x), where 0<a<1. Magenta printed answers fill reflection, -f(x), vertical translation, f(x+h), af(x) where a>1, vertical compression, af(x) where 0<a<1.}
\answer{reflection; (-f(x)); vertical translation; (f(x+h)); (af(x)), where (a>1); vertical compression; (af(x)), where (0<a<1).}
\end{te-addendum}
\begin{te-addendum}{Assess}{GeneralizeETP}
\textbf{Digital Differentiated Resources}
The Reteach to Build Understanding, Additional Practice, and Enrichment activities are available as digital assignments. These activities are automatically assigned when students complete the lesson quiz online and are automatically scored.
\textbf{Digital activity screenshot}
Solve the system of equations by graphing.
(y=3x+4)
(y=-2x+4)
Use the graphing tool to graph the system.
\placeholder{graph}{Digital exercise blank Cartesian grid with x/y axes and arrowheads, origin O, unit grid lines, numbered -10,-8,-6,-4,-2,2,4,6,8,10 on both axes. Window [-10,10] by [-10,10]; no functions or points plotted. Question Help menu overlays upper-right quadrant.}
Click to enlarge graph.
Click the graph, choose a tool in the palette and follow the instructions to create your graph.
\te{Interface labels: Exit; Question Help; Help Me Solve This; View an Example; Video; Textbook; Glossary; Math Tools; Print; 1 part remaining; Clear All; Check Answer; Review progress; Question 5 of 14; Go; Back; Next. Resource previews bear “enVision Algebra 2 • Teaching Resources,” SavvasRealize.com, and a copyright notice.}
\end{te-addendum}

\end{document}
